\scanpage{651}\mspage{\deleted{\unsure{38}}39}
pour $I$ quelconque dans Cat
\[
 \Phi_I\overset{\mathrm{def}}{=}\Phi\langle I\rangle\overset{\mathrm{def}}{=}\Psi(I)^*=v_!(\Phi(I_0))^{\sim},\tag*{(3.5.2)}
\]
où $\Psi=\Phi_*$, et où on suppose $\Phi$ Q-clos à g. i.e.
$\Phi=\Psi^*$. Dans le cas où on suppose
seulement $\Phi$ Q-saturé transfini, sans plus,
on aurait envie de poser
\[
 \Phi_I\overset{\mathrm{def}}{=}\text{Q-saturé transfini de }v_!(\Phi(I_0)),\tag*{(3.5.3)}
\]
et je présume que pour $I$ ordonné noethérien,
cela coïncide avec la définition d'Anderson.
Mais je vais me borner pour la suite au cas où
$C=\Phi$ est Q-clos à g.

Revenons alors au cas d'une structure à
cofibrations $(W,C)$ sur $M$, en supposant
$C$ et $TC=C\cap W$ Q-clos à gauche. Définissons
pour tout $I$ de Cat $C_I=C\langle I\rangle$ comme
en 3.5.2. On a alors, pour $u:I\unsure{-|\to}I'$ une
flèche de Cat
\[
 u_!(C\langle I\rangle)\subset C\langle I'\rangle,\tag*{(3.5.4)}
\]
mais il n'y a guère de raison d'espérer
avoir aussi (3.4.2), sans disposer de la
formule
\[
 C\langle I\rangle\cap W(I)\overset{?}{=}(TC)\langle I\rangle.\tag*{(3.5.5)}
\]
\margin{NB. On a l'inclusion $\supset$ ; car $v^*v_!TC(\unsure{I_0|I}\unsure{)}\subset TC(\unsure{I_1|I_0})$}
On a prouvé que cette formule est valable si $I$
est ordonné noethérien ($\underline{\overline{\mathrm{XVII}}}$ \inserted{th.~2.1} p.~26--27). Cela
% archive sheet number: 650
